\def\EditionMode{1}
\def\EditionLabel{Local math-prose proof}
\def\EditionStatus{Proposed wrappers; canonical source unchanged}
\documentclass[10pt,oneside]{book}
\usepackage{fontspec}
\setmainfont{Latin Modern Roman}
\setsansfont{Latin Modern Sans}
\usepackage{amsmath,amssymb,mathtools}
\usepackage{unicode-math}
\setmathfont{Latin Modern Math}
\usepackage{xcolor,tikz-cd,lua-ul,marginnote,pdflscape}
\usetikzlibrary{arrows.meta,decorations.pathmorphing,calc,shapes.geometric}
\tikzcdset{ampersand replacement=\&}
\tikzcdset{every label/.append style={text=black}}
\usepackage[a4paper,left=18mm,right=45mm,top=18mm,bottom=22mm,marginparwidth=28mm,marginparsep=7mm]{geometry}
\usepackage[unicode,colorlinks=true,urlcolor=black,linkcolor=black]{hyperref}
\usepackage{bookmark}
\setlength{\parindent}{0pt}
\setlength{\parskip}{4pt plus 1pt minus 1pt}
\setlength{\emergencystretch}{1.5em}
\raggedbottom
% Include final-line descenders inside the declared text-height rectangle.
\maxdepth=0pt
\definecolor{PSHand}{RGB}{25,58,132}
\definecolor{PSDelete}{RGB}{146,37,37}
\definecolor{PSTypedOld}{RGB}{95,95,95}
\ifnum\EditionMode=2\colorlet{blue}{black}\else\colorlet{blue}{PSHand}\fi
\RenewDocumentCommand{\textcircled}{m}{\tikz[baseline=(PSCircle.base)]\node[draw,circle,inner sep=.5pt,line width=.4pt](PSCircle){#1};}
\newif\ifPSInApparatus
\newcounter{psnote}
\newcounter{psvariant}
\newcounter{pslinkmark}
\newif\ifPSInLink
\newsavebox{\PSNoteBox}
\newdimen\PSNoteLimit
{\footnotesize\global\PSNoteLimit=3\baselineskip}
\newcommand{\CurrentScanURL}{}
\newcommand{\CurrentKey}{}
\newcommand{\CurrentScan}{}
\newcommand{\pssource}[3]{\psflushmargins\clearpage\gdef\CurrentKey{#1}\gdef\CurrentScan{#2}\gdef\CurrentScanURL{#3}\setcounter{psvariant}{0}\typeout{PS-SOURCE: #1}}
\newcommand{\psdivision}[2]{\pdfbookmark[0]{#1}{#2}\par\noindent{\large\bfseries #1\par}\medskip}
\newcommand{\pslinkmarker}[2]{\stepcounter{pslinkmark}\latelua{texio.write_nl('log','PS-LINK-SHIP|#1|\CurrentKey|\arabic{pslinkmark}|#2')}}
\newcommand{\pslink}[2]{\ifPSInLink\pslinkmarker{#1}{nested}#2\else\begingroup\PSInLinktrue\href{\CurrentScanURL}{\pslinkmarker{#1}{outer}#2}\endgroup\fi}
\newcommand{\scanpage}[1]{\ifnum\EditionMode=2\leavevmode\marginpar{\scriptsize\sffamily\pslink{scanpage}{scan #1}}\else\par\noindent{\small\sffamily\pslink{scanpage}{[scan #1]}}\par\medskip\fi}
\newcommand{\mspage}[1]{\ifnum\EditionMode=2\leavevmode\marginpar{\scriptsize\sffamily p.\,#1}\else{\small\sffamily[p.\,#1]}\fi}
\newcommand{\hand}[1]{{\ifnum\EditionMode=2\color{black}\else\color{PSHand}\fi #1}}
\let\over\relax
\NewDocumentCommand{\over}{O{} m m}{{\textcolor{PSTypedOld}{#2}\ensuremath{^{\!\nearrow}}\if\relax\detokenize{#1}\relax\hand{#3}\else #3\fi}}
\NewDocumentCommand{\inserted}{O{} m}{{\if\relax\detokenize{#1}\relax\hand{\ensuremath{\langle}#2\ensuremath{\rangle}}\else\ensuremath{\langle}#2\ensuremath{\rangle}\fi}}
\newcommand{\deleted}[1]{{\color{PSDelete}\ifmmode\text{[del.: }#1\text{]}\else[\textsf{del.:} #1]\fi}}
\newcommand{\xout}[1]{{\color{PSTypedOld}\ifmmode\text{[x: }#1\text{]}\else[\textsf{x:} #1]\fi}}
\let\PSMathUnderline\underline
\RenewDocumentCommand{\underline}{+m}{\ifmmode\PSMathUnderline{#1}\else\underLine{#1}\fi}
\NewDocumentCommand{\uline}{+m}{\underline{#1}}
\NewDocumentCommand{\uuline}{+m}{\ifmmode\PSMathUnderline{\PSMathUnderline{#1}}\else\underLine{\underLine{#1}}\fi}
\newunderlinetype\PSWaveLine{\leaders\hbox{\lower.7ex\hbox{\tikz[x=1pt,y=1pt,baseline=0pt]{\draw[line width=.3pt] (0,0) .. controls (1,1) and (2,1) .. (3,0) .. controls (4,-1) and (5,-1) .. (6,0);}}}}
\NewDocumentCommand{\uwave}{+m}{{\PSWaveLine #1}}
\newcommand{\unsure}[1]{\pslink{unsure}{[#1\ifmmode\,?\else\,?\fi]}}
\newcommand{\illegible}[1]{\ifmmode\text{\pslink{illegible}{[illeg.\ #1]}}\else\pslink{illegible}{[illeg.\ #1]}\fi}
\NewDocumentCommand{\diagram}{+m}{\ifmmode\vcenter{\hbox{#1}}\else\par\medskip{\centering #1\par}\medskip\psflushmargins\fi}
\newenvironment{diagramblock}{\par\begin{center}}{\end{center}\par}
\newcommand{\PSMathNotes}{}
\makeatletter
\NewDocumentCommand{\psplacednote}{m +m}{%
      \sbox{\PSNoteBox}{\parbox{\marginparwidth}{\footnotesize\raggedright\hand{\textsf{m#1.} #2}}}%
      \ifdim\dimexpr\ht\PSNoteBox+\dp\PSNoteBox\relax<\PSNoteLimit
        \leavevmode\marginpar{\usebox{\PSNoteBox}}%
      \else
        \par\nobreak\smallskip\hrule\nobreak\smallskip
        {\footnotesize\textsf{Marginal note m#1.}\quad\hand{#2}\par}
        \nobreak\smallskip\hrule\smallskip
      \fi}
\newcommand{\psflushmargins}{\ifmmode\else\ifinner\else\ifx\PSMathNotes\@empty\else\let\PSTempNotes\PSMathNotes\global\let\PSMathNotes\@empty\PSTempNotes\fi\fi\fi}
\NewDocumentCommand{\margin}{+m}{%
  \ifPSInApparatus\ifmmode\hand{\text{[marg.: }#1\text{]}}\else\hand{\textsf{[marg.: }#1\textsf{]}}\fi
  \else\stepcounter{psnote}%
    \ifmmode
      \text{\textsuperscript{m\arabic{psnote}}}%
      \begingroup\edef\PSQueue{\endgroup\noexpand\g@addto@macro\noexpand\PSMathNotes{\noexpand\psplacednote{\arabic{psnote}}{\unexpanded{$#1$}}}}\PSQueue
    \else
      \textsuperscript{m\arabic{psnote}}\psplacednote{\arabic{psnote}}{#1}%
    \fi
  \fi}
\AtBeginDocument{\everymath\expandafter{\the\everymath\aftergroup\psflushmargins}\everydisplay\expandafter{\the\everydisplay\aftergroup\psflushmargins}}
\makeatother
\newcommand{\psnumber}[2]{\expandafter\gdef\csname psnum@#1\endcsname{#2}}
\NewDocumentCommand{\var}{+m +m m m}{\ifnum\EditionMode=0 #1\else\ifnum\EditionMode=1\ifmmode\text{\textsuperscript{\tiny\csname psnum@#4\endcsname}}\else\textsuperscript{\tiny\csname psnum@#4\endcsname}\fi\fi #2\fi}
\newcommand{\psapparatusstart}{\par\bigskip\hrule\smallskip\PSInApparatustrue\begingroup\footnotesize\raggedright\textsf{Apparatus for scan \CurrentScan}\par}
\newcommand{\psapparatusend}{\par\endgroup\PSInApparatusfalse}
\NewDocumentCommand{\psentry}{m m m +m +m m}{\par\noindent\textbf{#2. #3.}\quad\textit{Typescript:} #4\quad\textit{Corrected:} \if\relax\detokenize{#5}\relax\textnormal{[omitted]}\else #5\fi\quad #6\par}
\newcommand{\psfrontmatter}{%
\hypersetup{pageanchor=false}\begin{titlepage}\raggedright
{\Large Alexander Grothendieck\par}\vspace{12mm}
{\Huge Pursuing Stacks\par}\medskip{\LARGE 1983\par}\vspace{12mm}
{\Large\EditionLabel\par}\vspace{8mm}
{\large\EditionStatus\par}\vfill
Université de Montpellier, Grothendieck archive, cotes 134-2--134-8.\par
Including the 26 source pages of R. Brown's translation at cote 134-2, sites 115--143.\par
\vspace*{2mm}\end{titlepage}\hypersetup{pageanchor=true}
\chapter*{Editorial statement}
AI-transcribed by Claude Opus 5.5 (Anthropic, effort high); checked against the scans, and this edition built, by GPT-6 Astra (OpenAI, ChatGPT Pro mode); prompts and records by Claude (Anthropic).

The witness consulted and not reproduced is the LaTeX edition of Mateo Carmona with the collaboration of Ulrik Buchholtz (arXiv 2111.01000 v2, 2021), extended from the 2015 edition by ``the scrivener'', CC0 1.0.

The translation in the appendix is R. Brown's, prepared with assistance from L. Breen and J-L. Loday, as his note says; the rights in it belong to its translator or his heirs.

The rights in the original work belong to its author or the author's representatives or heirs. Where we hold rights, in the AI-produced transcription, translation, corrections and typesetting, they are dedicated to the public domain under CC0 1.0. This edition is published for educational, non-profit purposes, and we claim no rights we do not hold. If you hold rights in this work and object to its publication here, please open a GitHub issue and we will take it down.

\section*{How to read this form}
\ifnum\EditionMode=0
The diplomatic form preserves typed and handwritten wording, cancellations, insertions, overwritten readings, underlining, doubts, and original page numbers. Blue marks handwriting; red labels handwritten cancellations; gray labels typewriter cancellations and underlying overwritten text. Angle brackets mark insertions; a superscript arrow separates an overwritten reading from its replacement. Typed insertions and typed overwriting remain plain. The disposition of long marginal notes is identified explicitly.
\else
The corrected text applies the author's insertions, replacements, cancellations, and recorded transpositions. It retains his English and notation, including consistent forms such as ``standart'', ``developped'', ``it's'' for ``its'', and ``A la''. An isolated slip is corrected only when specifically recorded and justified. Mathematical anomalies are preserved and discussed in the guide.
\ifnum\EditionMode=1
Blue continues to identify handwriting. A numbered call identifies each maximal editorial operation, with its full diplomatic and corrected readings and reason in the apparatus following the source page. Internal layers of a wholly replaced or deleted passage are preserved within that entry rather than duplicated as separate entries. These entries are generated from the same variants that govern the running text.
\else
This reading form omits the apparatus, cancelled material, superseded readings, insertion brackets, and color distinctions. It retains the resulting text, marginal notes, underlining, and unresolved doubts. Scan and source-page references appear in the margin.
\fi\fi

Every scan mark and remaining doubt links to its source page in the archive viewer. Doubts are editorial reading marks, not emendations. The source-page breaks identify the archival sequence; a source page may occupy more than one printed page. A partial working edition is explicitly labelled and contains only reviewed source pages.
\clearpage}

\begin{document}
\pssource{ps_134-5_001}{134-5/1}{https://grothendieck.umontpellier.fr/web/viewer.html?file=../134-5.pdf\#page=2}
\psnumber{ps_134-5_001-v0001}{1}
\psnumber{ps_134-5_001-v0002}{2}
\psnumber{ps_134-5_001-v0003}{3}
\psnumber{ps_134-5_001-v0004}{4}
\psnumber{ps_134-5_001-v0005}{5}
\psnumber{ps_134-5_001-v0006}{6}
\psnumber{ps_134-5_001-v0007}{7}
\psnumber{ps_134-5_001-v0008}{8}
\psnumber{ps_134-5_001-v0009}{9}
\psnumber{ps_134-5_001-v0010}{10}
\psnumber{ps_134-5_001-v0011}{11}
\psnumber{ps_134-5_001-v0012}{12}
\psnumber{ps_134-5_001-v0013}{13}
\psnumber{ps_134-5_001-v0014}{14}
\psnumber{ps_134-5_001-v0015}{15}
% Pursuing Stacks (cotes 134-2 to 134-8): archive cote 134-5, site page 1 (PDF page 2), typescript page 188.
% typescript. Draft compared: ps4.tex:1-63.
% AI-transcribed by Claude Opus 5.5 (claude-opus-5-5, adaptive thinking, effort high), Message Batches API, batch ps2 (msgbatch_01NVwr1Yuo8qJf5rykeGBXPK), 2026-11-12; reviewed against both supplied scan halves; see CHECK_DECISIONS.tsv.
\scanpage{134-5/1}
\mspage{188}
11.4.

6\var{\over{4}{7}}{\hand{7}}{writing over}{ps_134-5_001-v0001}.    A little more pondering and scribbling finally seems to show
that the \var{\deleted{real}}{}{deletion}{ps_134-5_001-v0002} key for an understanding of modelizers is'nt really the
notion of contractibility,  but rather the notion of aspheric objects
(besides,  of course,  the notion of weak equivalence).  At the same time
it appears that the notion of an \uline{aspheric} \uline{map} in (Cat),  more specifical-%
ly of a $\underline{W}$-aspheric``map''\var{\xout{xxxx}}{}{typewriter deletion}{ps_134-5_001-v0003} (i.e. functor between small categories) is
a lot more important than \var{\inserted{being just}}{\hand{being just}}{insertion}{ps_134-5_001-v0004} a highly expedient technical convenience,  as
it has been so far  -  it is indeed one of the basic notions of the theo-%
ry of modelizers we got into.  As a matter of fact,  I should have known
this for a number of weaks already, \var{\inserted{ever}}{\hand{ever}}{insertion}{ps_134-5_001-v0005} since I did some scribbling about
the plausible notion of ``morphism'' between test-categories (as well as
their weak and strong variants),  and readily convinced myself that the
natural ``morphisms'' here were nothing else but the aspheric functors bet-%
ween those categories.  I kind of forgot about this,  as it did'nt seem
too urgent to start moving \var{\inserted{around}}{\hand{around}}{insertion}{ps_134-5_001-v0006} the test category I was working with.  If I \var{\xout{xx}}{}{typewriter deletion}{ps_134-5_001-v0007}
had been a little more systematic in grinding through the usual functori-%
alities,  as soon as a significant notion (such as the various test noti-%
ons)  appears,  I presumably would have hit upon the crucial point about
modelizers and so-called ``asphericity structures'' a lot sooner,  without
going through the long-winded detour of homotopy structures\hand{,} and \var{\inserted{the still}}{\hand{the still}}{insertion}{ps_134-5_001-v0008} extreme-%
ly special types of test functors suggested by contractibility assumpti-%
ons.  However \var{\xout{it}}{}{typewriter deletion}{ps_134-5_001-v0009} I believe that most of the work I went through,  \var{\xout{xxx}}{}{typewriter deletion}{ps_134-5_001-v0010}
although irrelevant for the ``asphericity story'' itself,  will still be
useful,  especially when it comes to pinpointing the so-called ``canonic\unsure{\hand{al}}
modelizers,  whose modelizing structure is intrinsicly determined by the
category structure.

First thing now which we've to do is to have a closer look at the
meaning of asphericity for \var{\inserted{a}}{\hand{a}}{insertion}{ps_134-5_001-v0011} functor\var{\xout{s}}{}{typewriter deletion}{ps_134-5_001-v0012} between  small categories.  There is
no reason whatever to \var{\xout{restrict to}}{}{typewriter deletion}{ps_134-5_001-v0013} make any restrictions on these catego-%
ries besides smallness  (namely the cardinals of the sets of objects and
arrows being in the ``universe'' we are implicitly working in,  or even
more stringently still,  A and B to be in (Cat)  namely to be objects
of that universe).  Thus we will not assume A and B to be test categories
or the like.  \var{\xout{However,}}{}{typewriter deletion}{ps_134-5_001-v0014} We will be led to consider,  for \uline{any} small cate-%
gory A ,  the localization of A\^{} with respect to $\underline{W}$-equivalences,  which
I'll denote by $(\mathrm{Hot}^{\hand{\underline{W}}}_A)$ \hand{or simply $(\mathrm{Hot}_A)$:}
\[ \hand{(1)} \qquad \hand{(\mathrm{Hot}^{\underline{W}}_A)\text{ or}} \ (\mathrm{Hot}^{\var{\deleted{\illegible{1}}}{}{deletion}{ps_134-5_001-v0015}}_A) \ = \ \underline{W}_A^{-1}A\hat{} \quad . \]

\margin{1}

\psapparatusstart
\psentry{ps_134-5_001-v0001}{1}{writing over}{\over{4}{7}}{\hand{7}}{The superposed handwritten reading replaces the underlying reading.}
\psentry{ps_134-5_001-v0002}{2}{deletion}{\deleted{real}}{}{The author cancels this material by hand; it leaves the corrected running text.}
\psentry{ps_134-5_001-v0003}{3}{typewriter deletion}{\xout{xxxx}}{}{The typewriter strikes out this material; it leaves the corrected running text.}
\psentry{ps_134-5_001-v0004}{4}{insertion}{\inserted{being just}}{\hand{being just}}{The handwritten insertion is incorporated at its indicated position; insertion brackets leave the running text.}
\psentry{ps_134-5_001-v0005}{5}{insertion}{\inserted{ever}}{\hand{ever}}{The handwritten insertion is incorporated at its indicated position; insertion brackets leave the running text.}
\psentry{ps_134-5_001-v0006}{6}{insertion}{\inserted{around}}{\hand{around}}{The handwritten insertion is incorporated at its indicated position; insertion brackets leave the running text.}
\psentry{ps_134-5_001-v0007}{7}{typewriter deletion}{\xout{xx}}{}{The typewriter strikes out this material; it leaves the corrected running text.}
\psentry{ps_134-5_001-v0008}{8}{insertion}{\inserted{the still}}{\hand{the still}}{The handwritten insertion is incorporated at its indicated position; insertion brackets leave the running text.}
\psentry{ps_134-5_001-v0009}{9}{typewriter deletion}{\xout{it}}{}{The typewriter strikes out this material; it leaves the corrected running text.}
\psentry{ps_134-5_001-v0010}{10}{typewriter deletion}{\xout{xxx}}{}{The typewriter strikes out this material; it leaves the corrected running text.}
\psentry{ps_134-5_001-v0011}{11}{insertion}{\inserted{a}}{\hand{a}}{The handwritten insertion is incorporated at its indicated position; insertion brackets leave the running text.}
\psentry{ps_134-5_001-v0012}{12}{typewriter deletion}{\xout{s}}{}{The typewriter strikes out this material; it leaves the corrected running text.}
\psentry{ps_134-5_001-v0013}{13}{typewriter deletion}{\xout{restrict to}}{}{The typewriter strikes out this material; it leaves the corrected running text.}
\psentry{ps_134-5_001-v0014}{14}{typewriter deletion}{\xout{However,}}{}{The typewriter strikes out this material; it leaves the corrected running text.}
\psentry{ps_134-5_001-v0015}{15}{deletion}{$\deleted{\illegible{1}}$}{}{The author cancels this material by hand; it leaves the corrected running text.}
\psapparatusend

\end{document}
